\documentclass[9pt]{extarticle}

\usepackage[utf8]{inputenc}
\usepackage[a4paper, margin=1.2cm]{geometry}
\usepackage[colorlinks=true, urlcolor=blue]{hyperref}
\usepackage{enumitem}
\usepackage{changepage}

\pagestyle{empty}

\setlist[itemize]{label=\textbullet, noitemsep, topsep=0pt}
\setlength{\parindent}{0pt}

\begin{document}

\begin{center}
	{\huge\bfseries Paul Parisot}
	\\[4pt]
	{\Large Ingénieur Logiciel - C++, Performance \& Optimisation}
	\\
\end{center}

\par\vspace{-5pt}
\rule{\linewidth}{0.4pt}
\par\vspace{-5pt}

\begin{center}
	Github: \href{https://github.com/paulogarithm}{paulogarithm}
	\\
	\href{mailto:paul.parisot@epitech.eu}{paul.parisot@epitech.eu} - +33 6 42 36 01 57
	\\
\end{center}

\par\vspace{-10pt}
\rule{\linewidth}{0.4pt}
\par\vspace{7pt}

{\large\bfseries À propos}
\par\vspace{-8pt}
\rule{\linewidth}{0.4pt}
\par\vspace{5pt}

Élève-ingénieur en dernière année (Epitech, diplôme en 2027), à la recherche d'un {\bfseries stage de fin d'études}. Profil {\bfseries C++ et performance} : optimisation, gestion mémoire, SIMD, avec l'habitude de {\bfseries reprendre et corriger du code existant} (C++ legacy chez Forvia) et de mesurer les gains (temps par image).
\\
Contributions open-source sous CI/CD stricte (Gno.land), tests, Valgrind, aucune fuite mémoire.
\\
Bases en mathématiques et en informatique théorique (préparation de l'agrégation d'informatique). Un an de formation en management à McGill University.
\par\vspace{5pt}

{\large\bfseries Compétences}
\par\vspace{-8pt}
\rule{\linewidth}{0.4pt}
\par\vspace{5pt}

\hspace*{5pt}{\bfseries Langages}: C++ (C++23/26), C, Go, Python, Shell, Assembleur x86
\\
\hspace*{5pt}{\bfseries Performance}: optimisation, SIMD, alignement mémoire, multithreading, architecture orientée données (ECS), Linux
\\
\hspace*{5pt}{\bfseries Outils}: Git, Docker, Kubernetes, CI/CD, gdb, Valgrind
\\
\hspace*{5pt}{\bfseries Langues}: Français, Anglais (C1, compétences professionnelles)

\par\vspace{5pt}

{\large\bfseries Expérience Professionnelle}
\par\vspace{-8pt}
\rule{\linewidth}{0.4pt}
\par\vspace{5pt}

{\bfseries Gno.land - Ingénieur Logiciel Junior (Open-Source, Go)} \hfill avr. -- juil. 2025
\begin{itemize}[noitemsep, topsep=0pt]
	\item Contributions à un projet open-source : développement de {\bfseries 6+ packages core en Gnolang} (variante de Go pour smart contracts), dont une compatibilité EVM ; {\bfseries ~10k lignes de code} contribuées.
	\item Détection et correction de {\bfseries 6 bugs critiques} causant des crashs runtime.
	\item Rédaction de tests exhaustifs sous des politiques strictes de CI/CD et de revue de code.
\end{itemize}

\par\vspace{5pt}

{\bfseries Faurecia / Forvia - Stagiaire Ingénieur Logiciel} \hfill juil. -- déc. 2023
\begin{itemize}[noitemsep, topsep=0pt]
	\item Reprise et correction d'un ancien code {\bfseries C++} pour le faire communiquer avec les bons drivers matériels (LED) ; scripts Python de pilotage de hardware.
	\item Conception d'un {\bfseries moteur de calcul d'itinéraire GPS} et d'une interface de recherche de POI, intégrés à une application démo présentée au {\bfseries CES 2024}.
\end{itemize}

\par\vspace{5pt}

{\bfseries Epitech / IONIS - Assistant Technique \& Formateur} \hfill févr. 2024 -- mars 2025
\begin{itemize}[noitemsep, topsep=0pt]
	\item Enseignement à {\bfseries plus de 400 étudiants} en C, C++, asm et Haskell (dont multithreading) ; correction et évaluation des rendus.
\end{itemize}

\par\vspace{5pt}

{\large\bfseries Projets Personnels}
\par\vspace{-8pt}
\rule{\linewidth}{0.4pt}
\par\vspace{5pt}

\underline{\bfseries \href{https://github.com/paulogarithm/blorox}{BloRox} -- Moteur 3D temps réel en C++}
\par
\begin{adjustwidth}{5pt}{0pt}
    Architecture ECS orientée données : {\bfseries 100k entités animées}, $\sim$6 ms de moteur et $\sim$6 ms de rendu par image ; alignement mémoire et extensions SIMD.
\end{adjustwidth}
\par
{\bfseries Compétences :} C++23/26, ECS, optimisation de performance, SIMD

\par\vspace{8pt}

\underline{\bfseries \href{https://github.com/paulogarithm/cclass}{cclass} -- Moteur de mutation mémoire à l'exécution}
\par
\begin{adjustwidth}{5pt}{0pt}
    Framework en C pur réécrivant à l'exécution des instructions assembleur en mémoire ({\bfseries \texttt{mmap}}, flags de protection, alignement strict) pour contourner le surcoût du dispatch par vtable.
\end{adjustwidth}
\par
{\bfseries Compétences :} C, systèmes Linux, assembleur x86, mmap, sûreté mémoire

\par\vspace{8pt}

\underline{\bfseries Zappy -- Serveur TCP multi-clients en C}
\par
\begin{adjustwidth}{5pt}{0pt}
    Serveur TCP mono-thread à boucle événementielle ({\bfseries \texttt{select}}), codé seul, avec gestion des clients par liste chaînée de connexions. \\
    Buffer borné par client et mise en quarantaine des clients qui le saturent, pour éviter tout épuisement mémoire ; aucune erreur ni fuite sous Valgrind.
\end{adjustwidth}
\par
{\bfseries Compétences :} C, sockets TCP, select, gestion mémoire, Valgrind, tests unitaires et fonctionnels

\par\vspace{8pt}

\underline{\bfseries Sledge -- Moteur de migration, synchronisation et orchestration multi-cloud (Go)}
\par
\begin{adjustwidth}{5pt}{0pt}
    Outil de migration multi-cloud en {\bfseries Go} pour des flux de données à haut volume, via une couche d'abstraction cloud personnalisée ({\ttfamily provicol}) supportant AWS S3 et GCP ; gestion de structures de données complexes et d'interactions avec des bases de données ({\bfseries intégrité des données}).
\end{adjustwidth}
\par
{\bfseries Compétences :} Go, SQL, backend haut volume

\par\vspace{8pt}



{\large\bfseries Formation}
\par\vspace{-8pt}
\rule{\linewidth}{0.4pt}
\par\vspace{5pt}

{\bfseries Epitech, Paris} \hfill 2022 -- 2027 \\
Diplôme d'ingénieur en technologies de l'information (diplôme de niveau Master) \\
Certification Niveau 7 - Expert en ingénierie logicielle

\vspace{6pt}

{\bfseries McGill University, Montréal} \hfill 2025 -- 2026 \\
Certificat en Management -- GPA : 3.57/4.0 \\
Cours suivis : Comptabilité, Managerial Economics \& Analysis, Business, Mathématiques, Gestion de projet, Marketing

\end{document}
